All deviations were logged when they occurred, in the time-stamped file \texttt{protocol/DEVIATIONS.md}; they are reproduced here verbatim except for formatting.
\begin{itemize}
\item \textbf{D1 (P0 gate) — TimesFM-2.5 XReg batching error, fixed and rerun.} First gate run (kernel lta-gate-k1) passed all 500 series to one \texttt{forecast\_with\_covariates} call. The XReg library (\texttt{timesfm/utils/xreg\_lib.py}, \texttt{fit}) solves a single least-squares problem over the stacked rows of all series in a call, i.e. covariate coefficients are shared across the batch. With independent series whose effect signs are random, the shared coefficient is near zero: all covariate arms gave GM-MASE 0.6583-0.6590. This is an implementation error in our wrapper, not a protocol choice. Fix: one call per series (per-series in-context regression, the documented intent). TimesFM-2.5 gate rerun as kernel lta-gate-k3-timesfm; the first run is retained in results/gate\_k1 and reported as the "batch-shared XReg" variant.
\item \textbf{D2 (synthetic generator) — \texttt{null} family added after the gate.} A \texttt{null} family (beta = 0, covariate uninformative) was added to \texttt{lta/synth.py} for LTA training and as a do-no-harm test. It was not part of the gate cells; gate results are unaffected (seeds depend on family code).
\item \textbf{D3 (P2 rivers) — pipeline dry run touched 61 test windows.} Before the full run, a CPU dry run of the main kernel (Chronos-2-small, LTA trained 3 steps, every \textasciitilde{}150th window per split: 61 train / 61 val / 61 test windows) and of the fine-tuning kernel (20 test windows) were executed to check the code path. No design choice (arms, hyper-parameters, hypotheses) was changed after it. Its numbers are not reported.
\item \textbf{D4 (P2 rivers) — LTA river fine-tuning never improved validation; LTA-ft = LTA-zs.} In all runs (main seed 0, seeds 1-2, ablations), every fine-tuning checkpoint (lr 1e-3, 150-step evaluation) had worse val GM-MASE than the step-0 zero-shot aligner, so pre-registered early stopping returned the step-0 weights: LTA-ft and LTA-lopo are numerically identical to LTA-zs. Local diagnosis (12 CPU steps): the river quantile loss is heavy-tailed (batch losses 0.15-1.07) and the aligner sharpens its temperature and shortens its lead (known horizon fraction 0.70 $\rightarrow$ 0.61). Reported as is; H2/H3 are evaluated with the identical weights.
\item \textbf{D5 (P2 transfer) — TabPFN-TS evaluated on every 4th test window.} TabPFN-TS needs \textasciitilde{}1.2 s per window on a T4; all arms on 9,046 windows would exceed the Kaggle session limit. A fixed deterministic subsample (window index \% 4 == 0; 2,262 windows, all pairs) is used for every TabPFN-TS arm. Chronos-2, TimesFM-2.5 and Moirai-1.1 use all test windows.
\item \textbf{D6 (Chronos-2 LoRA baseline) — first attempt crashed (torchao/peft version clash); rerun with torchao removed.} 
\item \textbf{D7 (transfer) — float16 overflow in saved aligned covariates.} The main kernel stored aligned covariates as float16, which overflows above 65,504 cfs (Mississippi, Missouri, Ohio). The in-kernel Chronos-2 evaluation used float32 and is unaffected. For transfer, the aligned covariates were regenerated in float32 on CPU from the saved seed-0 weights (lta\_zs.pt); vs the GPU float16 copy: median relative difference 1.9e-4, mask disagreement 0.014\% of positions (threshold ties). To keep inputs identical across backbones, Chronos-2 is re-evaluated on the regenerated arrays as well (arm \texttt{lta\_zs}, transfer run).
\item \textbf{D8 (transfer) — Chronos-2 float32 re-evaluation and Moirai transfer run on local CPU.} The Kaggle transfer kernels for Chronos-2 and Moirai-1.1 stayed queued (account GPU-session limits). Both were run locally on CPU with tools/local\_transfer.py: same test windows, same arm construction (lta/arms.py), same float32 aligned arrays (D7), same model checkpoints (amazon/chronos-2; Salesforce/moirai-1.1-R-base). Only the device differs; Moirai uses 100 samples and the median as in the GPU wrapper.
\item \textbf{D9 (AirLead) — post-hoc specificity check (not pre-registered).} Because the pre-registered random-lag control beat every aligned arm on AirLead test, an exploratory check was run on every 12th VALIDATION window (305 windows, Chronos-2): fixed lags 24/48/96/168 h, CCF restricted to [24,168] h, and an i.i.d. noise channel known over context and horizon. Reported as post hoc (appendix table of the post-hoc AirLead check); it does not change any pre-registered outcome.
\item \textbf{D10 (statistics) — Wilcoxon tie tolerance.} Paired synthetic MASE differences below 1e-6 (float32 rounding of identical backbone inputs) are set to exact zeros before the Wilcoxon test. Without this, p-values for families with many ties differed across processes (e.g. linear: 0.873 vs 0.899) because of last-bit differences in CSV round-trips. Conclusions are unchanged.
\item \textbf{D11 (transfer) — TabPFN-TS run split into four disjoint parts.} The single TabPFN-TS transfer kernel stayed queued. The identical window set (every 4th test window, D5) was split deterministically into four disjoint interleaved parts (idx[part::4]) run as separate kernels (two on T4, two on P100). Arms, weights and inputs are unchanged; results are concatenated. Redundant queued kernels for Chronos-2 and Moirai (run locally, D8) were deleted unrun. The original single kernel eventually ran; it covers all 2,262 windows and reproduces the three completed parts exactly (8,480 overlapping rows, zero difference). The full run is used; the remaining queued part was deleted unrun.
\item \textbf{D12 (TimesFM-2.5 wrapper) — no-covariate arm read the backcast instead of the forecast.} TimesFM-2.5 was compiled with return\_backcast=True (required for XReg). With this setting the library prepends the backcast to the point forecast, and our no-covariate path sliced the FIRST H values (padded context) instead of the LAST H. Only the TimesFM \texttt{uni} arm is affected (gate, RiverLead transfer, AirLead); all covariate arms go through forecast\_with\_covariates, which slices the horizon correctly (verified in the library source). Found because the river \texttt{uni} GM-MASE (2.91) was implausible. Fix: take the last H values; the \texttt{uni} arm was rerun with the same checkpoint on the identical task sets: locally on CPU for the gate (tools/timesfm\_uni\_fix.py, which also recomputes the old slicing; old GM-MASE 0.837 reproduced exactly, corrected 0.659), and on Kaggle GPU for RiverLead and AirLead test together with the random-lag control (D13). Gate rule G1 for TimesFM is re-evaluated with the corrected reference.
\item \textbf{D13 (TimesFM-2.5) — post-hoc random-lag control on RiverLead and AirLead.} The pre-registered TimesFM-2.5 arms did not include the random-lag control. Because the control changed the interpretation of AirLead for Chronos-2, it was added post hoc for TimesFM-2.5 on the RiverLead and AirLead test splits (same seeds and lag range [H, 168] as for Chronos-2), together with the corrected no-covariate arm (D12). Recorded before these runs finished; reported as post hoc. Comparisons between LTA and classical arms are unaffected (both add one aligned channel to the raw covariate; top-3 adds three).
\item \textbf{D14 (P5-A) — synthetic aligned covariates regenerated on CPU; pulse family numerically sensitive.} For the synthetic transfer runs the three frozen aligners were applied on CPU in float32. Re-scoring a subset with Chronos-2 on CPU reproduced the GPU per-series MASE to a median relative difference of 4e-7, except in the sparse pulse family, where tiny numerical differences in the near-zero aligned values are amplified by the backbone's scaling (25 pulse series: GM local/GPU 1.038, sd of log ratio 0.126; time-varying-lag series identical). All transfer comparisons are within one backbone and use the same CPU arrays for every LTA arm.
\item \textbf{D15 (P5-A) — dry run touched 26 synthetic test series.} A CPU dry run of the synthetic transfer kernel (TimesFM-2.5, 2 series per cell, 5 arms) checked the code path. Nothing was changed after it and its numbers are not reported.
\item \textbf{D16 (P5-B) — binned classical lag: fallback when no window passes the correlation threshold.} fit\_binned\_lag keeps only training windows whose CCF peak is >= 0.3. On EnergyLead some pairs have few such windows (GB$\rightarrow$NL: 4 of 132 in the subsample used by the estimator), and in the subsampled dry run a pair had none, which crashed the estimator. Fallback added: if no window qualifies, the peaks of all windows are used. This case never occurred on RiverLead or AirLead, so their results are unchanged. Found in the dry run (D17), before any energy forecast was scored.
\item \textbf{D17 (P5-B) — EnergyLead dry run.} A CPU dry run (Chronos-2-small, \textasciitilde{}150 windows per split, val only, 4 arms) checked the code path. Its numbers are not reported and nothing was changed after it except D16.
\item \textbf{D18 (P5-B) — TimesFM-2.5 EnergyLead runs every classical arm.} To avoid waiting for the Chronos-2 validation selection, TimesFM-2.5 is run on the test split with all classical arms (ccf, ccf\_pw, gcc\_phat, top3, global, binned) in addition to uni, raw, randlag and the three LTA seeds. The hypothesis tests use only the arm selected on Chronos-2 validation, as pre-registered; other arms are descriptive.
\item \textbf{D19 (P5-D) — CoRA implementation details and a post-hoc no-covariate control.} (a) The paper zero-initialises the projection and "the MLP". Zero-initialising every MLP layer gives zero gradients for all parameters, so we follow the adaLN-Zero convention the paper cites (DiT): projection and the MLP output layer zero-initialised, MLP first layer default-initialised. Verified: at initialisation the CoRA median forecast equals frozen univariate Chronos-2 (max abs difference 0.031 on values \textasciitilde{}6.8e4, float32). (b) Because the adapter's bias terms can learn a covariate-independent recalibration of Chronos-2 on each benchmark, a post-hoc control \texttt{cora\_nocov} (same adapter and training, covariate embedding set to zero) is added. It separates the effect of the covariate from that of in-domain adaptation.
\item \textbf{D20 (P5-D) — CoRA dry run.} A CPU dry run (Chronos-2-small, \textasciitilde{}60 windows per split, 6 steps) checked the code path on all three domains. Its numbers are not reported and nothing was changed after it.
\item \textbf{D21 (P5-C) — post-hoc log-response LightGBM on RiverLead.} The pre-registered supervised baselines model y\_{t+h} - y\_t in discharge units. Because discharge is heavy-tailed and hydrological models usually work in log space, a post-hoc variant (\texttt{lgbm\_full\_log}) uses log(y) for all target features, log(x) for covariate features and log y\_{t+h} - log y\_t as response (non-positive discharge, 35 hourly values at the Ohio downstream gauge, is floored at 0.1\% of the pair's training median), with otherwise identical settings (full history, validation early stopping). Added after the pre-registered RiverLead supervised results were seen; reported as post hoc. The pre-registered arms are unchanged.
\item \textbf{D22 (P6) — TrafficLead pair screening outcome; descriptive short-lag subset defined before any forecast.} The pre-registered selection (train split only) gave 10 pairs. Five have train global lags of 3-29 steps (15 min to 2.4 h: 101S\_400743 29, 85S\_400528 4, 85S\_400209 8, 85N\_400336 3, 85N\_400148 8); five have 103-112 steps (8.6-9.3 h: 680S\_407184, 101N\_402359, 680S\_407185, 101N\_401908, 87N\_404452), which is implausible for congestion propagation over 6-12 miles and likely reflects the daily commute pattern. Train CCF peaks are weak (0.080-0.108). The primary analysis keeps all 10 pairs as pre-registered. A descriptive 'short-lag' subset (the five pairs with train lag < 60 steps) is also reported. It was defined from training data only, before any validation or test forecast, and is labelled post hoc.
\item \textbf{D23 (P6) — TrafficLead dry run.} A CPU dry run of the evaluation kernel (Chronos-2-small, \textasciitilde{}60 windows per split, val only, 6 arms) checked the code path, as did a dry run of the CoRA kernel (6 steps, \textasciitilde{}60 windows). Their numbers are not reported and nothing was changed after them.
\item \textbf{D24 (post hoc, all domains) — training-split lead diagnostic.} After the TrafficLead Chronos-2 result (no aligned arm beat the raw covariate although train lags were >= 3 steps), the train global lag and the mean train CCF peak value at that lag (same estimator as the screening, log- differences) were computed for every pair of all four real benchmarks and related to the per-pair test ratio of the validation-selected classical arm (and LTA seed 0) to the raw covariate with Chronos-2. Descriptive only; computed after the test results of RiverLead, AirLead and EnergyLead were known.
\item \textbf{D25 (P6) — TrafficLead: imputed sensors found; quality-filtered subset.} The supervised baselines gave implausible LightGBM GM-MASE of 0.005-0.02 on two pairs. Inspection showed that their target sensors (401908, 404452) repeat their value exactly one week earlier at 100\% of validation and test steps and \textasciitilde{}92\% of training steps, i.e. they carry imputed, not measured, speeds. Two further sensors exceed 50\% exact weekly repeats in one split (target 402359: 92\% train; covariate 400665: 56\% val); all other sensors are at 1-7\%. The protocol did not anticipate imputed series. Found after the Chronos-2 test results were known. We report both the pre-registered 10 pairs and a quality-filtered set of 6 pairs that excludes any pair with a sensor above 50\% exact weekly repeats in any split (excluded: 101N\_401908, 87N\_404452, 101N\_402359, 680S\_407184). The Chronos-2 conclusions are the same on both sets. Supervised and cross-domain comparisons use the filtered set, because a model can reproduce imputed values exactly.
\end{itemize}